% Source: physical PDF pp. 453--459 (printed pp. 430--436). Coverage:
% complete Section 23.2, from its heading through immediately before
% Section 23.3; Eqs. (23.2.1)--(23.2.9), nine unnumbered displays,
% no citations, figures, tables, or footnotes.
\subsection{Homotopy Groups}
\label{sec:23.2}

We learned in the previous section to classify field configurations, at
which the Hamiltonian or other functionals are finite, in correspondence
with the elements of appropriate homotopy groups. But we have not yet
explained in what sense the the homotopy groups are \emph{groups}, nor have
we given any physical significance to the group structure. As we shall see,
there is a natural definition of the multiplication rule for elements of
homotopy groups, according to which two extended configurations of fields
forming a manifold $\mathcal M$ in $d$ dimensions, that belong to different
elements $c_1$ and $c_2$ of $\pi_d(\mathcal M)$, can only fuse continuously
to form a configuration belonging to the element $c_1\cdot c_2$ of
$\pi_d(\mathcal M)$.

We will begin by defining the first homotopy group $\pi_1(\mathcal M)$ of
an arbitrary manifold $\mathcal M$, also known as the \emph{fundamental
group of the manifold}. As we have seen, the existence of a non-trivial
$\pi_1(G/H)$ for some coset space $G/H$ is the condition for the
topological stability of a vortex line in three dimensions (or a monopole
in two dimensions). After considering $\pi_1(\mathcal M)$, we will then
move on to more general homotopy groups.

A connected manifold $\mathcal M$ is said to be multiply connected if
there is some closed curve of points $p(z)$ on the manifold, parameterized
by a single variable $z$ with $0\leq z\leq1$ and $p(0)=p(1)$, which cannot
be contracted to a point by a continuous deformation. Since on a connected
manifold we can always continuously deform any such closed curve so that
any one point on the curve is anywhere we like on the manifold, we may
restrict our attention only to curves for which
$p(0)=p(1)=p_0$, where $p_0$ is any fixed point of the manifold, known as
the \emph{base point}. Two such closed curves $p_1(z)$ and $p_2(z)$ are
said to be homotopically equivalent if they can be deformed into each
other, that is, if there is a continuous function $p(z,t)$ for
$0\leq t\leq1$, with
\[
 p(z,0)=p_1(z),\qquad p(z,1)=p_2(z);
\]
\[
 p(0,t)=p(1,t)=p_0.
\]
The relation of homotopic equivalence is an equivalence relation in the
sense of being symmetric, reflexive, and transitive, so it divides the
space of closed curves on the manifold into equivalence classes: two
closed curves are in the same class if and only if they are homotopically
equivalent. The set of these equivalence classes is known as the first
homotopy group of the manifold, $\pi_1(\mathcal M)$.

To define the multiplication rule for $\pi_1(\mathcal M)$, choose a
standard curve $p[z;c]$ that starts and ends at the base point $p_0$ for
each equivalence class $c$ in $\pi_1(\mathcal M)$. For any two equivalence
classes $c_1$ and $c_2$, define the `product' $c_1\cdot c_2$ as the
equivalence class containing the curve $p[z,c_1,c_2]$ that starts at
$p_0$, follows $p[z,c_1]$ back to $p_0$, and then follows $p[z,c_2]$ back
again to $p_0$. Formally, we take
\[
 p[z,c_1,c_2]\equiv
 \begin{cases}
  p[2z,c_1],&0\leq z\leq\tfrac12,\\
  p[2z-1,c_2],&\tfrac12\leq z\leq1.
 \end{cases}
\]

We must now show that multiplication defined in this way satisfies the
conditions for a group. First, let us check that this multiplication is
associative. For this purpose, note that $(c_1\cdot c_2)\cdot c_3$ is the
equivalence class containing a curve $p[z,c_1\cdot c_2,c_3]$ that goes
along the standard curve $p[z,c_1\cdot c_2]$ from the base point and back
again, and then along the standard curve $p[z,c_3]$ from the base point
and back again, while $c_1\cdot(c_2\cdot c_3)$ is the equivalence class
containing a curve $p[z,c_1,c_2\cdot c_3]$ that goes along the standard
curve $p[z,c_1]$ from the base point and back again, and then along the
standard curve $p[z,c_2\cdot c_3]$ from the base point and back again. By
definition, the curve $p[z,c_1\cdot c_2]$ may be deformed into a curve that
goes along $p[z,c_1]$ from the base point and back again and then along
$p[z,c_2]$ from the base point and back again, while
$p[z,c_2\cdot c_3]$ may be deformed into a curve that goes along
$p[z,c_2]$ from the base point and back again and then along $p[z,c_3]$
from the base point and back again. Hence both curves
$p[z,c_1\cdot c_2,c_3]$ and $p[z,c_1,c_2\cdot c_3]$ may be deformed into
the curve that goes along $p[z,c_1]$ from the base point and back again,
then along $p[z,c_2]$ from the base point and back again, and finally
along $p[z,c_3]$ from the base point and back again, and hence they may
be deformed into each other, showing that
\[
 (c_1\cdot c_2)\cdot c_3=c_1\cdot(c_2\cdot c_3).
\]

The unit element $e$ of $\pi_1(\mathcal M)$ is defined as the equivalence
class containing the curve $p[z,e]=p_0$ that stays at the base point. To
check that $e\cdot c=c$, note that
\[
 p[z,e,c]=
 \begin{cases}
  p_0,&0\leq z\leq\tfrac12,\\
  p[2z-1,c],&\tfrac12\leq z\leq1.
 \end{cases}
\]
But this curve can be continuously deformed into $p[z,c]$ by taking
\[
 p(z,t)=
 \begin{cases}
  p_0,&0\leq z\leq t/2,\\
  p[(2z-t)/(2-t),c],&t/2\leq z\leq1,
 \end{cases}
\]
which is the same as $p[z,e,c]$ for $t=1$ and the same as $p[z,c]$ for
$t=0$. The product $e\cdot c$ is the equivalence class containing
$p[z,e,c]$, which we now see is the same as the equivalence class
containing $p[z,c]$, which is just $c$. The proof that $c\cdot e=c$ is
similar.

The `inverse' $c^{-1}$ of the equivalence class $c$ is the equivalence
class that contains a curve $p^{-1}[z,c]$ which goes around the same path
as the standard curve $p[z,c]$ but in the opposite direction; that is,
\[
 p^{-1}[z,c]=p[1-z,c].
\]
This is not necessarily the same as the `standard' curve $p[z,c^{-1}]$,
but by the definition of the equivalence class $c^{-1}$, the two curves
may be deformed into each other. To check that $c^{-1}\cdot c=e$, note
that by deforming $p[z,c^{-1}]$ into $p^{-1}[z,c]$, the curve
$p[z,c^{-1},c]$ may be deformed into
\[
 p[z,c^{-1},c]\longrightarrow
 \begin{cases}
  p[1-2z,c],&0\leq z\leq\tfrac12,\\
  p[2z-1,c],&\tfrac12\leq z\leq1.
 \end{cases}
\]
But this curve can be continuously deformed into $p[z,e]=p_0$ by taking
\[
 p(z,t)=
 \begin{cases}
  p[1-2tz,c],&0\leq z\leq\tfrac12,\\
  p[2tz+1-2t,c],&\tfrac12\leq z\leq1,
 \end{cases}
\]
which is the same as $p[z,c^{-1},c]$ for $t=1$ and the same as $p[z,e]$
for $t=0$. The product $c^{-1}\cdot c$ is the equivalence class containing
$p[z,c^{-1},c]$, which we now see is the same as the equivalence class
containing $p[z,e]$, which is just $e$. The proof that
$c\cdot c^{-1}=e$ is similar. The existence of a unit element and inverses
shows that these equivalence classes form a group.

The classic example of a manifold $\mathcal M$ with a non-trivial first
homotopy group is the circle itself, $\mathcal M=S_1$. This may be
parameterized by an angle $\theta$ with $\theta=0$ and
$\theta=2\pi n$ (with $n$ any positive or negative integer) identified as
the same point. The homotopy groups consist of classes of functions
$\theta(z)$ for $0\leq z\leq1$ that begin at some base point,
$\theta(0)=\theta_0$, and end at the same base point,
$\theta(1)=\theta_0+2\pi n$. Two such functions may be continuously
deformed into each other if and only if they have the same value of $n$,
so $\pi_1(S_1)$ consists of a denumerably infinite number of classes $c_n$
labelled by the positive or negative integer $n$. Furthermore, the
`product' of two classes $c_n$ and $c_m$ consists of curves that go $n$
times around the circle from the base point and back again, and then $m$
times around the circle from the base point and back again, so here
multiplication is just addition:
\begin{equation}
 c_n\cdot c_m=c_{n+m}.
 \tag{23.2.1}\label{eq:23.2.1}
\end{equation}
and hence
\begin{equation}
 \pi_1(S_1)=\mathbb Z,
 \tag{23.2.2}\label{eq:23.2.2}
\end{equation}
where $\mathbb Z$ is the group of addition of positive and negative
integers. As an immediate physical application, we note that when a gauge
group $SO(2)$ is completely spontaneously broken, the coset space is just
$SO(2)$ itself, which has the topology of a circle, so in this case there
is an infinite number of types of topologically stable vortex line,
characterized by a positive or negative integer $n$. For instance, this is
the case in a type II superconductor, where as we saw in
\hyperref[sec:21.6]{Section 21.6} the $U(1)$ of electromagnetic gauge
invariance is spontaneously broken to a discrete subgroup $\mathbb Z_2$.

More generally, all spheres $S_d$ with $d>1$ are simply connected, which
means that they have a trivial first homotopy group, a statement
conventionally expressed as
\begin{equation}
 \pi_1(S_d)=0\qquad\text{for }d>1.
 \tag{23.2.3}\label{eq:23.2.3}
\end{equation}
Only a few of the more familiar Lie groups are multiply connected:
\begin{equation}
 \pi_1(G)=
 \begin{cases}
  \mathbb Z,&G=U(k),\quad k\geq1,\\
  \mathbb Z_2,&G=SO(k),\quad k\geq3,\\
  0,&G=\operatorname{Spin}(k),\quad k\geq3,\\
  0,&G=SU(k),\quad k\geq2,\\
  0,&G=USp(2k),\quad k\geq1,\\
  0,&G=G_2,F_4,E_6,E_7,E_8.
 \end{cases}
 \tag{23.2.4}\label{eq:23.2.4}
\end{equation}
Here $\mathbb Z_2$ is the group with two elements $\mathbf 1$ and
$-\mathbf 1$, with group multiplication defined as ordinary
multiplication, and $\operatorname{Spin}(n)$ is the simply connected
covering group of $SO(n)$. (As we saw in
\hyperref[sec:2.7]{Section 2.7}, $\operatorname{Spin}(3)$ is the same as
$SU(2)$.) Also, for a direct product of two manifolds $\mathcal M$ and
$\mathcal M'$, the fundamental group is
\begin{equation}
 \pi_1(\mathcal M\cdot\mathcal M')
 =\pi_1(\mathcal M)\cdot\pi_1(\mathcal M').
 \tag{23.2.5}\label{eq:23.2.5}
\end{equation}

We can appreciate the physical significance of the group structure of
$\pi_1(\mathcal M)$ by asking what happens when two distant parallel
vortex lines in three dimensions are brought together. When the vortex
lines are sufficiently far apart their fields do not interact, so the
configuration can be described by specifying the classes $c'$ and $c''$
in $\pi_1(G/H)$ to which each belongs. The class $c$ to which the whole
configuration belongs is determined by the behavior of the fields on a
very large circle surrounding both vortex lines. By a continuous
deformation we can distort this very large circle into two large circles,
one surrounding each vortex line, which intersect at a point midway
between them. As we go around this closed curve in two-dimensional space
we first trace out a curve in $G/H$ which consists of one of the closed
curves in the class $c'$, and then one of the closed curves in the class
$c''$, just as in the definition of the product of the classes. We
conclude then that the whole configuration is in a class
$c=c'\cdot c''$, and so the two vortex lines can only fuse together to
form a vortex line of this class. In particular, they can only annihilate
if $c''=c'^{-1}$.

For instance, when $\pi_1(G/H)=\mathbb Z$ (as in superconductivity, where
$G/H=SO(2)/\mathbb Z_2$), in three dimensions two vortex lines of classes
$c_n$ and $c_m$ can fuse together to form a vortex line of class $c_{n+m}$,
so they can only annihilate if $n=-m$. On the other hand, when
$\pi_1(G/H)=\mathbb Z_2$ (as when $G=SO(N)$ with $N\geq3$ and $H$ is
trivial or discrete) there is only one kind of vortex line in three
dimensions, corresponding to the element $-\mathbf 1$ of $\mathbb Z_2$,
and since $(-\mathbf 1)^2=\mathbf 1$, any two can annihilate.

Now let us consider the general homotopy group $\pi_k(\mathcal M)$. This
is much like $\pi_1(\mathcal M)$, except that instead of considering
mappings of the circle $S_1$ into a manifold $\mathcal M$, we consider
mappings of the $k$-sphere $S_k$ (the surface of a $(k+1)$-dimensional
ball) into $\mathcal M$, again with one point of $S_k$ always mapped into
the same `base point' $p_0$ of $\mathcal M$. Two such mappings are
equivalent if one can be continuously deformed into the other, keeping
the same point of $S_k$ always mapped into the base point. The $k$th
homotopy group $\pi_k(\mathcal M)$ has elements consisting of equivalence
classes of these mappings.

It is often convenient to picture the $d$-sphere $S_d$ as the interior of
a $d$-dimensional hypercube, with all points on the boundary identified as
a single point. For instance, we have already seen that the circle $S_1$
can be treated as the interval $0\leq\theta\leq2\pi$, with points
$\theta=0$ and $\theta=2\pi$ identified. Similarly we can make a map of
an $S_2$ like the earth's surface by cutting out the south pole and
spreading out the resulting sheet on the unit square
$0\leq z_1\leq1$, $0\leq z_2\leq1$. Continuous mappings of this square
into $\mathcal M$ must take all points on the boundary into the same point
of $\mathcal M$, because all points on the boundary are the same point,
the south pole. In general, two mappings
$p(z_1,\ldots,z_d)$ and $p'(z_1,\ldots,z_d)$ of $S_d$ into $\mathcal M$
are homotopically equivalent if one can be continuously deformed into the
other, while keeping $p$ on the boundary of the hypercube equal to the
base point $p_0$.

As before, for each equivalence class $c$ we choose a standard mapping
$p(z_1,\ldots,z_d;c)$. The product of $c_1$ and $c_2$ is defined as the
equivalence class that contains the mapping
\begin{equation}
 p(z_1,z_2,\ldots,z_d;c_1,c_2)=
 \begin{cases}
  p(2z_1,z_2,\ldots,z_d;c_1),&0\leq z_1\leq\tfrac12,\\
  p(2z_1-1,z_2,\ldots,z_d;c_2),&\tfrac12\leq z_1\leq1.
 \end{cases}
 \tag{23.2.6}\label{eq:23.2.6}
\end{equation}
The unit element $e$ is defined as the equivalence class that contains the
mapping with $p=p_0$ for all $z$, and the inverse $c^{-1}$ of $c$ is
defined as the equivalence class that contains the mapping with
\begin{equation}
 p^{-1}[z_1,z_2,\ldots,z_N;c]
 =p[1-z_1,z_2,\ldots,z_N;c].
 \tag{23.2.7}\label{eq:23.2.7}
\end{equation}
In the same way as for $\pi_1$, it can be shown that this multiplication is
associative, and that $e\cdot c=c\cdot e=c$ and
$c^{-1}\cdot c=c\cdot c^{-1}=e$. All $\pi_n(\mathcal M)$ for $n\geq2$
are Abelian. (There are manifolds $\mathcal M$ for which
$\pi_1(\mathcal M)$ is non-Abelian, such as the plane with two or more
points removed.)

In any case in which $\pi_k(\mathcal M)=\mathbb Z$, there must be a
one-to-one mapping of the $k$-sphere $S_k$ into a $k$-sphere $S'_k$ in
$\mathcal M$, which corresponds to the element `one' of $\mathbb Z$ (not
the unit element, which is zero). The element $\nu$ of $\mathbb Z$ with
$\nu=2,3,\ldots$ corresponds to the mapping of $S_k$ into the same
$k$-sphere $S'_k$ in $\mathcal M$, which covers $S'_k$ $\nu$ times, with
the Jacobian of the transformation $S_k\to S'_k$ positive. The element
$\nu$ of $\mathbb Z$ with $\nu=-1,-2,\ldots$ corresponds to the mapping
$S_k\to S'_k$, which covers $S'_k$ $|\nu|$ times, with the Jacobian of
the transformation $S_k\to S'_k$ negative.

For instance, we saw in the previous section that magnetic monopoles arise
when a simply connected group $G$ is broken to the $U(1)$ of
electromagnetism. In this case the appendix shows that
\begin{equation}
 \pi_2(G/U(1))=\pi_1(U(1))=\mathbb Z,
 \tag{23.2.8}\label{eq:23.2.8}
\end{equation}
so a magnetic monopole carries an integer-valued quantum number $\nu$,
which as shown in \hyperref[sec:23.3]{Section 23.3} is proportional to the
magnetic charge. This quantum number gives the number of times that a
two-sphere of large radius surrounding the monopole is mapped into a
two-sphere in the manifold $G/U(1)$ of Goldstone boson fields (with the
relative orientation of the two two-spheres being the same or opposite
for $\nu$ positive or negative, respectively) and is therefore known as
the \emph{winding number}. The structure of the group $\mathbb Z$ shows
that this quantum number is conserved, in the sense that a monopole of
quantum number $\nu$ can fuse with a monopole of quantum number $\nu'$
only to form a monopole with quantum number $\nu+\nu'$.

If the unbroken subgroup is $SO(n)$ with $n\geq3$, then according to
\hyperref[app:23.B]{Appendix B of this chapter}
\begin{equation}
 \pi_2(G/SO(n))=\pi_1(SO(n))=\mathbb Z_2.
 \tag{23.2.9}\label{eq:23.2.9}
\end{equation}
In this case there is just one sort of `monopole,' corresponding to the
element $-\mathbf 1$ of $\mathbb Z_2$, which can annihilate only in
pairs. It is important to distinguish between this case and that in which
$SO(n)$ is replaced with its simply connected covering group
$\operatorname{Spin}(n)$, for which there are no monoples. We will come
back to this point at the end of the next section.

Another example: we saw in the previous section that the skyrmions in
quantum chromodynamics with $n$ light quarks correspond to elements of
$\pi_3(SU(n))$, which according to the appendix is $\mathbb Z$. This
these skyrmions carry a conserved integer-valued quantum number $\nu$,
which perhaps may be identified with baryon number. Similarly, recall
from the previous section that the instantons in a gauge theory based on
a simple gauge group $G$ correspond to the elements of $\pi_3(G)$, which
according to the appendix is $\mathbb Z$, so these instantons like
skyrmions carry an integer-valued quantum number $\nu$, also known as the
winding number. In \hyperref[sec:23.5]{Section 23.5} we shall see how to
express this quantum number as a local functional of the gauge field.
